\cleardoubleevenpage
\begin{colorbigpic}
	[\basicscale]
	{3ii_verse}
	[Hang there, my verse, in witness of my love]
\end{colorbigpic}
\Scene{2}[The Forest of Arden]
\enter{\textsc{Orlando,} with a paper}


\begin{verse_speech}[Orlando]
Hang there, my verse, in witness of my love.\\
   And thou, thrice-crownèd  queen of night, survey\\
With thy chaste eye, from thy pale sphere above,\\
   Thy huntress' name that my full life doth sway.\\
O Rosalind, these trees shall be my books,\\
   And in their barks my thoughts I'll character,\\
That every eye which in this forest looks\\
   Shall see thy virtue witnessed everywhere.\\
Run, run, Orlando, carve on every tree\\
The fair, the chaste, and unexpressive she.
\end{verse_speech}

\exit{\textsc{Orlando}}

\enter{\textsc{Corin} and \textsc{Touchstone}}


\begin{prose_speech}[Corin]
And how like you this shepherd's life, Master Touchstone?
\end{prose_speech}

\begin{prose_speech}[Touchstone]
Truly, shepherd, in respect of itself, it is a good life; but in respect that it is a shepherd's life, it is naught. In respect that it is solitary, I like it very well; but in respect that it is private, it is a very vile life. Now in respect it is in the fields, it pleaseth me well; but in respect it is not in the court, it is tedious. As it is a spare life, look you, it fits my humour well; but as there is no more plenty in it, it goes much against my stomach. Hast any philosophy in thee, shepherd?
\end{prose_speech}

\begin{prose_speech}[Corin]
No more but that I know the more one sickens, the worse at ease he is, and that he that wants money, means, and content is without three good friends; that the property of rain is to wet, and fire to burn; that good pasture makes fat sheep; and that a great cause of the night is lack of the sun; that he that hath learned no wit by nature nor art may complain of good breeding or comes of a very dull kindred.
\end{prose_speech}

\begin{colorbigpic}
	[\basicscale]
	{3ii_shepherd}
	[Wast ever in court, shepherd?]
\end{colorbigpic}

\begin{prose_speech}[Touchstone]
Such a one is a natural philosopher. Wast ever in court, shepherd?
\end{prose_speech}

\begin{prose_speech}[Corin]
No, truly.
\end{prose_speech}

\begin{prose_speech}[Touchstone]
Then thou art damned.
\end{prose_speech}

\begin{prose_speech}[Corin]
Nay, I hope.
\end{prose_speech}

\begin{prose_speech}[Touchstone]
Truly, thou art damned, like an ill-roasted egg, all on one side.
\end{prose_speech}

\begin{prose_speech}[Corin]
For not being at court? Your reason.
\end{prose_speech}

\begin{prose_speech}[Touchstone]
Why, if thou never wast at court, thou never saw'st good manners; if thou never saw'st good manners, then thy manners must be wicked, and wickedness is sin, and sin is damnation. Thou art in a parlous state, shepherd.
\end{prose_speech}

\begin{prose_speech}[Corin]
Not a whit, Touchstone. Those that are good manners at the court are as ridiculous in the country as the behaviour of the country is most mockable at the court. You told me you salute not at the court but you kiss your hands. That courtesy would be uncleanly if courtiers were shepherds.
\end{prose_speech}

\begin{prose_speech}[Touchstone]
Instance, briefly. Come, instance.
\end{prose_speech}

\begin{prose_speech}[Corin]
Why, we are still handling our ewes, and their fells, you know, are greasy.
\end{prose_speech}

\begin{prose_speech}[Touchstone]
Why, do not your courtier's hands sweat? And is not the grease of a mutton as wholesome as the sweat of a man? Shallow, shallow. A better instance, I say. Come.
\end{prose_speech}

\begin{prose_speech}[Corin]
Besides, our hands are hard.
\end{prose_speech}

\begin{prose_speech}[Touchstone]
Your lips will feel them the sooner. Shallow again. A more sounder instance. Come.
\end{prose_speech}

\begin{prose_speech}[Corin]
And they are often tarred over with the surgery of our sheep; and would you have us kiss tar? The courtier's hands are perfumed with civet.
\end{prose_speech}

\begin{prose_speech}[Touchstone]
Most shallow man. Thou worms' meat in respect of a good piece of flesh, indeed. Learn of the wise and perpend: civet is of a baser birth than tar, the very uncleanly flux of a cat. Mend the instance, shepherd.
\end{prose_speech}

\begin{prose_speech}[Corin]
You have too courtly a wit for me. I'll rest.
\end{prose_speech}

\begin{prose_speech}[Touchstone]
Wilt thou rest damned? God help thee, shallow man. God make incision in thee; thou art raw.
\end{prose_speech}

\begin{prose_speech}[Corin]
Sir, I am a true labourer. I earn that I eat, get that I wear, owe no man hate, envy no man's happiness, glad of other men's good, content with my harm, and the greatest of my pride is to see my ewes graze and my lambs suck.
\end{prose_speech}

\begin{prose_speech}[Touchstone]
That is another simple sin in you, to bring the ewes and the rams together and to offer to get your living by the copulation of cattle; to be bawd to a bell-wether and to betray a she-lamb of a twelvemonth to a crooked-pated old cuckoldly ram, out of all reasonable match. If thou be'st not damned for this, the devil himself will have no shepherds. I cannot see else how thou shouldst 'scape.
\end{prose_speech}

\enter{\textsc{Rosalind,} as \textsc{Ganymede}}


\begin{prose_speech}[Corin]
Here comes young Master Ganymede, my new mistress's brother.
\end{prose_speech}

\begin{prose_speech}[Rosalind as Ganymede]
\stage{Reading a paper}
%\vspace{-1em}
	\begin{verse}\itshape
	\begin{altverse}
	From the east to western Ind\\
	No jewel is like Rosalind.\\
	Her worth being mounted on the wind,\\
	Through all the world bears Rosalind.\\
	All the pictures fairest lined\\
	Are but black to Rosalind.\\
	Let no face be kept in mind\\
	But the fair of Rosalind.
	\end{altverse}
\end{verse}
\end{prose_speech}

\begin{prose_speech}[Touchstone]
I'll rhyme you so eight years together, dinners and suppers and sleeping hours excepted. It is the right butter-women's rank to market.
\end{prose_speech}

\begin{prose_speech}[Rosalind as Ganymede]
  Out, fool.
\end{prose_speech}

\begin{prose_speech}[Touchstone]
For a taste:
% \vspace{-1em}
	\begin{verse}\itshape
	\begin{altverse}
	If a hart do lack a hind,\\
	Let him seek out Rosalind.\\
	If the cat will after kind,\\
	So be sure will Rosalind.\\
	Wintered garments must be lined;\\
	So must slender Rosalind.\\
	They that reap must sheaf and bind;\\
	Then to cart with Rosalind.\\
	Sweetest nut hath sourest rind;\\
	Such a nut is Rosalind.\\
	He that sweetest rose will find\\
	Must find love's prick, and Rosalind.
		\end{altverse}
\end{verse}
This is the very false gallop of verses. Why do you infect yourself with them?
\end{prose_speech}

\begin{prose_speech}[Rosalind as Ganymede] 
Peace, you dull fool. I found them on a tree.
\end{prose_speech}

\begin{prose_speech}[Touchstone]
Truly, the tree yields bad fruit.
\end{prose_speech}

\begin{prose_speech}[Rosalind as Ganymede] 
I'll graft it with you, and then I shall graft it with a medlar. Then it will be the earliest fruit i' th' country, for you'll be rotten ere you be half ripe, and that's the right virtue of the medlar.
\end{prose_speech}

\begin{prose_speech}[Touchstone]
You have said, but whether wisely or no, let the forest judge.
\end{prose_speech}

\enter{\textsc{Celia,} as \textsc{Aliena,} with a writing}


\begin{prose_speech}[Rosalind as Ganymede] 
Peace. Here comes my sister reading. Stand aside.
\end{prose_speech}

\begin{letter}
	\clearpage
\end{letter}

\begin{verse_speech}[Celia as Aliena]
\stage{Reads}
% \vspace{-1em}
	\begin{verse}\itshape
	\begin{altverse}
	Why should this a desert be?\\
	   For it is unpeopled? No.\\
	Tongues I'll hang on every tree\\
	   That shall civil sayings show.\\
	Some how brief the life of man\\
	   Runs his erring pilgrimage,\\
	That the stretching of a span\\
	   Buckles in his sum of age;\\
	Some of violated vows\\
	   'Twixt the souls of friend and friend.\\
	But upon the fairest boughs,\\
	   Or at every sentence' end,\\
	Will I <Rosalinda> write,\\
	   Teaching all that read to know\\
	The quintessence of every sprite\\
	   Heaven would in little show.\\
	Therefore heaven nature charged\\
	   That one body should be filled\\
	With all graces wide-enlarged.\\
	   Nature presently distilled\\
	Helen's cheek, but not her heart,\\
	   Cleopatra's majesty,\\
	Atalanta's better part,\\
	   Sad Lucretia's modesty.\\
	Thus Rosalind of many parts\\
	   By heavenly synod was devised\\
	Of many faces, eyes, and hearts\\
	   To have the touches dearest prized.\\
	Heaven would that she these gifts should have\\
	And I to live and die her slave.
		\end{altverse}
\end{verse}
\end{verse_speech}

\begin{prose_speech}[Rosalind as Ganymede] O most gentle Jupiter, what tedious homily of love have you wearied your parishioners withal, and never cried <Have patience, good people!>
\end{prose_speech}

\begin{prose_speech}[Celia as Aliena]
How now?—Back, friends. Shepherd, go off a little.—Go with him, sirrah.
\end{prose_speech}

\begin{prose_speech}[Touchstone]
Come, shepherd, let us make an honourable retreat, though not with bag and baggage, yet with scrip and scrippage.
\end{prose_speech}

\exit{\textsc{Touchstone} and \textsc{Corin}}

\begin{prose_speech}[Celia]
Didst thou hear these verses?
\end{prose_speech}

\begin{prose_speech}[Rosalind]
O yes, I heard them all, and more too, for some of them had in them more feet than the verses would bear.
\end{prose_speech}

\begin{prose_speech}[Celia]
That's no matter. The feet might bear the verses.
\end{prose_speech}

\begin{prose_speech}[Rosalind]
Ay, but the feet were lame and could not bear themselves without the verse, and therefore stood lamely in the verse.
\end{prose_speech}

\begin{prose_speech}[Celia]
But didst thou hear without wondering how thy name should be hanged and carved upon these trees?
\end{prose_speech}

\begin{prose_speech}[Rosalind]
I was seven of the nine days out of the wonder before you came, for look here what I found on a palm tree. 
\stage{She shows the paper she read.}
I was never so berhymed since Pythagoras' time that I was an Irish rat, which I can hardly remember.
\end{prose_speech}

\begin{prose_speech}[Celia]
Trow you who hath done this?
\end{prose_speech}

\begin{prose_speech}[Rosalind]
Is it a man?
\end{prose_speech}

\begin{prose_speech}[Celia]
And a chain, that you once wore, about his neck. Change you colour?
\end{prose_speech}

\begin{prose_speech}[Rosalind]
I prithee, who?
\end{prose_speech}

\begin{prose_speech}[Celia]
O Lord, Lord, it is a hard matter for friends to meet, but mountains may be removed with earthquakes and so encounter.
\end{prose_speech}

\begin{prose_speech}[Rosalind]
Nay, but who is it?
\end{prose_speech}

\begin{prose_speech}[Celia]
Is it possible?
\end{prose_speech}

\begin{prose_speech}[Rosalind]
Nay, I prithee now, with most petitionary vehemence, tell me who it is.
\end{prose_speech}

\begin{prose_speech}[Celia]
O wonderful, wonderful, and most wonderful wonderful, and yet again wonderful, and after that out of all whooping!
\end{prose_speech}

\begin{prose_speech}[Rosalind]
Good my complexion, dost thou think though I am caparisoned like a man, I have a doublet and hose in my disposition? One inch of delay more is a South Sea of discovery. I prithee, tell me who is it quickly, and speak apace. I would thou couldst stammer, that thou might'st pour this concealed man out of thy mouth as wine comes out of a narrow-mouthed bottle—either too much at once, or none at all. I prithee take the cork out of thy mouth, that I may drink thy tidings.
\end{prose_speech}

\begin{prose_speech}[Celia]
So you may put a man in your belly.
\end{prose_speech}

\begin{prose_speech}[Rosalind]
Is he of God's making? What manner of man? Is his head worth a hat, or his chin worth a beard?
\end{prose_speech}

\begin{prose_speech}[Celia]
Nay, he hath but a little beard.
\end{prose_speech}

\begin{prose_speech}[Rosalind]
Why, God will send more, if the man will be thankful. Let me stay the growth of his beard, if thou delay me not the knowledge of his chin.
\end{prose_speech}

\begin{prose_speech}[Celia]
It is young Orlando, that tripped up the wrestler's heels and your heart both in an instant.
\end{prose_speech}

\begin{prose_speech}[Rosalind]
Nay, but the devil take mocking. Speak sad brow and true maid.
\end{prose_speech}

\begin{prose_speech}[Celia]
I' faith, coz, 'tis he.
\end{prose_speech}

\begin{prose_speech}[Rosalind]
Orlando?
\end{prose_speech}

\begin{prose_speech}[Celia]
Orlando.
\end{prose_speech}

\begin{prose_speech}[Rosalind]
Alas the day, what shall I do with my doublet and hose? What did he when thou saw'st him? What said he? How looked he? Wherein went he? What makes he here? Did he ask for me? Where remains he? How parted he with thee? And when shalt thou see him again? Answer me in one word.
\end{prose_speech}

\begin{colorbigpic}
	[\basicscale]
	{3ii_doublet}
	[What shall I do with my doublet and hose?]
\end{colorbigpic}

\begin{prose_speech}[Celia]
You must borrow me Gargantua's mouth first. 'Tis a word too great for any mouth of this age's size. To say ay and no to these particulars is more than to answer in a catechism.
\end{prose_speech}

\begin{prose_speech}[Rosalind]
But doth he know that I am in this forest and in man's apparel? Looks he as freshly as he did the day he wrestled?
\end{prose_speech}

\begin{prose_speech}[Celia]
It is as easy to count atomies as to resolve the propositions of a lover. But take a taste of my finding him, and relish it with good observance. I found him under a tree like a dropped acorn.
\end{prose_speech}

\begin{prose_speech}[Rosalind]
It may well be called Jove's tree when it drops forth such fruit.
\end{prose_speech}

\begin{prose_speech}[Celia]
Give me audience, good madam.
\end{prose_speech}

\begin{prose_speech}[Rosalind]
Proceed.
\end{prose_speech}

\begin{prose_speech}[Celia]
There lay he, stretched along like a wounded knight.
\end{prose_speech}

\begin{colorbigpic}
	[\basicscale]
	{3ii_wounded}
	[There lay he, stretched along like a wounded knight]
\end{colorbigpic}

\begin{prose_speech}[Rosalind]
Though it be pity to see such a sight, it well becomes the ground.
\end{prose_speech}

\begin{prose_speech}[Celia]
Cry <holla> to thy tongue, I prithee. It curvets unseasonably. He was furnished like a hunter.
\end{prose_speech}

\begin{prose_speech}[Rosalind]
O, ominous! He comes to kill my heart.
\end{prose_speech}

\begin{prose_speech}[Celia]
I would sing my song without a burden. Thou bring'st me out of tune.
\end{prose_speech}

\begin{prose_speech}[Rosalind]
Do you not know I am a woman? When I think, I must speak. Sweet, say on.
\end{prose_speech}

\begin{prose_speech}[Celia]
You bring me out.

\enter{\textsc{Orlando} and \textsc{Jaques}}

Soft, comes he not here?
\end{prose_speech}

\begin{prose_speech}[Rosalind]
'Tis he. Slink by, and note him.
\end{prose_speech}

\stage{\textsc{Rosalind} and \textsc{Celia} step aside.}

\begin{colorbigpic}
	[\basicscale]
	{3ii_slink}
	['Tis he. Slink by, and note him.]
\end{colorbigpic}

\begin{prose_speech}[Jaques]
\stage{To \textsc{Orlando}}
I thank you for your company, but, good faith, I had as lief have been myself alone.
\end{prose_speech}

\begin{prose_speech}[Orlando]
And so had I, but yet, for fashion sake, I thank you too for your society.
\end{prose_speech}

\begin{prose_speech}[Jaques]
God be wi' you. Let's meet as little as we can.
\end{prose_speech}

\begin{prose_speech}[Orlando]
I do desire we may be better strangers.
\end{prose_speech}

\begin{prose_speech}[Jaques]
I pray you mar no more trees with writing love songs in their barks.
\end{prose_speech}

\begin{prose_speech}[Orlando]
I pray you mar no more of my verses with reading them ill-favouredly.
\end{prose_speech}

\begin{prose_speech}[Jaques]
Rosalind is your love's name?
\end{prose_speech}

\begin{prose_speech}[Orlando]
Yes, just.
\end{prose_speech}

\begin{prose_speech}[Jaques]
I do not like her name.
\end{prose_speech}

\begin{prose_speech}[Orlando]
There was no thought of pleasing you when she was christened.
\end{prose_speech}

\begin{prose_speech}[Jaques]
What stature is she of?
\end{prose_speech}

\begin{prose_speech}[Orlando]
Just as high as my heart.
\end{prose_speech}

\begin{prose_speech}[Jaques]
You are full of pretty answers. Have you not been acquainted with goldsmiths' wives and conned them out of rings?
\end{prose_speech}

\begin{prose_speech}[Orlando]
Not so. But I answer you right painted cloth, from whence you have studied your questions.
\end{prose_speech}

\begin{prose_speech}[Jaques]
You have a nimble wit. I think 'twas made of Atalanta's heels. Will you sit down with me? And we two will rail against our mistress the world and all our misery.
\end{prose_speech}

\begin{prose_speech}[Orlando]
I will chide no breather in the world but myself, against whom I know most faults.
\end{prose_speech}

\begin{prose_speech}[Jaques]
The worst fault you have is to be in love.
\end{prose_speech}

\begin{prose_speech}[Orlando]
'Tis a fault I will not change for your best virtue. I am weary of you.
\end{prose_speech}

\begin{prose_speech}[Jaques]
By my troth, I was seeking for a fool when I found you.
\end{prose_speech}

\begin{prose_speech}[Orlando]
He is drowned in the brook. Look but in, and you shall see him.
\end{prose_speech}

\begin{prose_speech}[Jaques]
There I shall see mine own figure.
\end{prose_speech}

\begin{prose_speech}[Orlando]
Which I take to be either a fool or a cipher.
\end{prose_speech}

\begin{prose_speech}[Jaques]
I'll tarry no longer with you. Farewell, good Signior Love.
\end{prose_speech}

\begin{prose_speech}[Orlando]
I am glad of your departure. Adieu, good Monsieur Melancholy.
\end{prose_speech}

\exit{\textsc{Jaques}}

\begin{prose_speech}[Rosalind]
\asideto{Celia}{I will speak to him like a saucy lackey, and under that habit play the knave with him.}
\stage{As \textsc{Ganymede}}  Do you hear, forester?
\end{prose_speech}

\begin{colorbigpic}
	[\basicscale]
	{3ii_saucy}
	[I will speak to him like a saucy lackey]
\end{colorbigpic}

\begin{prose_speech}[Orlando]
Very well. What would you?
\end{prose_speech}

\begin{prose_speech}[Rosalind as Ganymede] I pray you, what is 't o'clock?
\end{prose_speech}

\begin{prose_speech}[Orlando]
You should ask me what time o' day. There's no clock in the forest.
\end{prose_speech}

\begin{prose_speech}[Rosalind as Ganymede] Then there is no true lover in the forest; else sighing every minute and groaning every hour would detect the lazy foot of time as well as a clock.
\end{prose_speech}

\begin{prose_speech}[Orlando]
And why not the swift foot of time? Had not that been as proper?
\end{prose_speech}

\begin{prose_speech}[Rosalind as Ganymede] By no means, sir. Time travels in divers paces with divers persons. I'll tell you who time ambles withal, who time trots withal, who time gallops withal, and who he stands still withal.
\end{prose_speech}

\begin{prose_speech}[Orlando]
I prithee, who doth he trot withal?
\end{prose_speech}

\begin{prose_speech}[Rosalind as Ganymede] Marry, he trots hard with a young maid between the contract of her marriage and the day it is solemnized. If the interim be but a se'nnight, time's pace is so hard that it seems the length of seven year.
\end{prose_speech}

\begin{prose_speech}[Orlando]
Who ambles time withal?
\end{prose_speech}

\begin{prose_speech}[Rosalind as Ganymede] With a priest that lacks Latin and a rich man that hath not the gout, for the one sleeps easily because he cannot study, and the other lives merrily because he feels no pain—the one lacking the burden of lean and wasteful learning, the other knowing no burden of heavy tedious penury. These time ambles withal.
\end{prose_speech}

\begin{prose_speech}[Orlando]
Who doth he gallop withal?
\end{prose_speech}

\begin{prose_speech}[Rosalind as Ganymede] With a thief to the gallows, for though he go as softly as foot can fall, he thinks himself too soon there.
\end{prose_speech}

\begin{prose_speech}[Orlando]
Who stays it still withal?
\end{prose_speech}

\begin{prose_speech}[Rosalind as Ganymede] With lawyers in the vacation, for they sleep between term and term, and then they perceive not how time moves.
\end{prose_speech}

\begin{prose_speech}[Orlando]
Where dwell you, pretty youth?
\end{prose_speech}

\begin{prose_speech}[Rosalind as Ganymede] With this shepherdess, my sister, here in the skirts of the forest, like fringe upon a petticoat.
\end{prose_speech}

\begin{prose_speech}[Orlando]
Are you native of this place?
\end{prose_speech}

\begin{prose_speech}[Rosalind as Ganymede] As the cony that you see dwell where she is kindled.
\end{prose_speech}

\begin{prose_speech}[Orlando]
Your accent is something finer than you could purchase in so removed a dwelling.
\end{prose_speech}

\begin{prose_speech}[Rosalind as Ganymede] I have been told so of many. But indeed an old religious uncle of mine taught me to speak, who was in his youth an inland man, one that knew courtship too well, for there he fell in love. I have heard him read many lectures against it, and I thank God I am not a woman, to be touched with so many giddy offences as he hath generally taxed their whole sex withal.
\end{prose_speech}

\begin{prose_speech}[Orlando]
Can you remember any of the principal evils that he laid to the charge of women?
\end{prose_speech}

\begin{prose_speech}[Rosalind as Ganymede] There were none principal. They were all like one another as halfpence are, every one fault seeming monstrous till his fellow fault came to match it.
\end{prose_speech}

\begin{prose_speech}[Orlando]
I prithee recount some of them.
\end{prose_speech}

\begin{prose_speech}[Rosalind as Ganymede] No, I will not cast away my physic but on those that are sick. There is a man haunts the forest that abuses our young plants with carving <Rosalind> on their barks, hangs odes upon hawthorns and elegies on brambles, all, forsooth, deifying the name of Rosalind. If I could meet that fancy-monger, I would give him some good counsel, for he seems to have the quotidian of love upon him.
\end{prose_speech}

\begin{prose_speech}[Orlando]
I am he that is so love-shaked. I pray you tell me your remedy.
\end{prose_speech}

\begin{prose_speech}[Rosalind as Ganymede] There is none of my uncle's marks upon you. He taught me how to know a man in love, in which cage of rushes I am sure you are not prisoner.
\end{prose_speech}

\begin{prose_speech}[Orlando]
What were his marks?
\end{prose_speech}

\begin{prose_speech}[Rosalind as Ganymede] A lean cheek, which you have not; a blue eye and sunken, which you have not; an unquestionable spirit, which you have not; a beard neglected, which you have not—but I pardon you for that, for simply your having in beard is a younger brother's revenue. Then your hose should be ungartered, your bonnet unbanded, your sleeve unbuttoned, your shoe untied, and everything about you demonstrating a careless desolation. But you are no such man. You are rather point-device in your accouterments, as loving yourself than seeming the lover of any other.
\end{prose_speech}

\begin{prose_speech}[Orlando]
Fair youth, I would I could make thee believe I love.
\end{prose_speech}

\begin{prose_speech}[Rosalind as Ganymede] Me believe it? You may as soon make her that you love believe it, which I warrant she is apter to do than to confess she does. That is one of the points in the which women still give the lie to their consciences. But, in good sooth, are you he that hangs the verses on the trees wherein Rosalind is so admired?
\end{prose_speech}

\begin{prose_speech}[Orlando]
I swear to thee, youth, by the white hand of Rosalind, I am that he, that unfortunate he.
\end{prose_speech}

\begin{prose_speech}[Rosalind as Ganymede] But are you so much in love as your rhymes speak?
\end{prose_speech}

\begin{prose_speech}[Orlando]
Neither rhyme nor reason can express how much.
\end{prose_speech}

\begin{prose_speech}[Rosalind as Ganymede] Love is merely a madness, and, I tell you, deserves as well a dark house and a whip as madmen do; and the reason why they are not so punished and cured is that the lunacy is so ordinary that the whippers are in love too. Yet I profess curing it by counsel.
\end{prose_speech}

\begin{prose_speech}[Orlando]
Did you ever cure any so?
\end{prose_speech}

\begin{prose_speech}[Rosalind as Ganymede] Yes, one, and in this manner. He was to imagine me his love, his mistress, and I set him every day to woo me; at which time would I, being but a moonish youth, grieve, be effeminate, changeable, longing and liking, proud, fantastical, apish, shallow, inconstant, full of tears, full of smiles; for every passion something, and for no passion truly anything, as boys and women are, for the most part, cattle of this colour; would now like him, now loathe him; then entertain him, then forswear him; now weep for him, then spit at him, that I drave my suitor from his mad humour of love to a living humour of madness, which was to forswear the full stream of the world and to live in a nook merely monastic. And thus I cured him, and this way will I take upon me to wash your liver as clean as a sound sheep's heart, that there shall not be one spot of love in 't.
\end{prose_speech}

\begin{prose_speech}[Orlando]
I would not be cured, youth.
\end{prose_speech}

\begin{prose_speech}[Rosalind as Ganymede] I would cure you if you would but call me Rosalind and come every day to my cote and woo me.
\end{prose_speech}

\begin{prose_speech}[Orlando]
Now, by the faith of my love, I will. Tell me where it is.
\end{prose_speech}

\begin{prose_speech}[Rosalind as Ganymede] Go with me to it, and I'll show it you; and by the way you shall tell me where in the forest you live. Will you go?
\end{prose_speech}

\begin{prose_speech}[Orlando]
With all my heart, good youth.
\end{prose_speech}

\begin{prose_speech}[Rosalind as Ganymede] Nay, you must call me Rosalind.—Come, sister, will you go?
\end{prose_speech}

\exeunt{}